% ===== Wrap-up (~4 min, then questions) =====
\partaccent{butter}\notopic

\begin{frame}{Այսօրվա 5 գաղափարը}
  \begin{enumerate}
    \item<1-> Անհատականությունը \textbf{ուղեղի} աշխատանք է (Գեյջ), բայց այն չափելը շատ դժվար է
    \item<2-> «Մեծ հնգյակը» գիտության լավագույն մոդելն է, իսկ MBTI-ն ու օնլայն թեստերը \textbf{զվարճանք} են, ոչ թե գիտություն
    \item<3-> Զայրույթը կարող է օգտակար լինել. իջեցնում է կորտիզոլն ու մղում է գործելու։ Բայց \textbf{չափն} է կարևոր, ու «թափելը» չի օգնում
    \item<4-> Խոստացված պարգևը կարող է \textbf{սպանել} ներքին մոտիվացիան, իսկ ինքնավարությունը՝ ուժեղացնել այն
    \item<5-> Հումորը խախտված ու լուծված սպասում է, ու այն \textbf{սոցիալական} է. ուրիշների հետ մոտ 30 անգամ ավելի հաճախ ենք ծիծաղում
  \end{enumerate}
\end{frame}

\begin{frame}{Փորձիր տանը}
  \begin{columns}[T]
    \begin{column}{0.31\textwidth}
      \begin{card}[colback=mint!30!cardwhite]
        \centering\bigemoji[1cm]{masks}\par\smallskip
        \raggedright Խնդրիր երեք ընկերոջ գնահատել քեզ Մեծ հնգյակի հինգ գծով (1-5)։ Համընկա՞վ քո գնահատականի հետ
      \end{card}
    \end{column}
    \begin{column}{0.31\textwidth}
      \begin{card}[colback=brainpink!30!cardwhite]
        \centering\bigemoji[1cm]{steam}\par\smallskip
        \raggedright Երբ հաջորդ անգամ անմեղի վրա բարկանաս, հաշվիր մինչև տասը ու մտածիր՝ ինչն էր կուտակվել օրվա ընթացքում
      \end{card}
    \end{column}
    \begin{column}{0.31\textwidth}
      \begin{card}[colback=lavender!35!cardwhite]
        \centering\bigemoji[1cm]{speech}\par\smallskip
        \raggedright Պատմիր ծխնելույզի կատակը ընկերոջդ ու հետևիր արձագանքին։ Նենսի Բելի փորձը՝ տանը
      \end{card}
    \end{column}
  \end{columns}
\end{frame}

\begin{frame}{Հաջորդ վեբինարը}
  \begin{columns}[c]
    \begin{column}{0.6\textwidth}
      \chip[mint]{hug}{7. Խմբային գրկախառնություն}\par\smallskip
      Ինչպես են ուրիշ մարդիկ ազդում ուղեղի վրա\par\medskip
      \begin{itemize}
        \item Ինչու է դժվար թաքցնել այն, ինչ իրականում մտածում ես
        \item Ինչպես ենք կառավարում ուրիշներին, ու ինչպես են մեզ կառավարում
        \item Ինչու է բաժանումն այդքան ցավոտ
        \item Ինչ է անում ուղեղը խմբում, ու ինչու ենք երբեմն ուրիշների հետ վատ վարվում
      \end{itemize}
    \end{column}
    \begin{column}{0.34\textwidth}
      \centering\mascot[2.6cm]{happy}\par\medskip
      {\small\color{plum!70!cream}Հիշու՞մ ես տապալված կատակը։ Բըրնեթն ասում էր, որ դա հասկանալու համար առանձին գլուխ է պետք։ Ահա և այն։}
    \end{column}
  \end{columns}
\end{frame}

\begin{frame}[plain]
  \begin{tikzpicture}[remember picture, overlay]
    \fill[lavender!60!cream] (current page.south west) rectangle (current page.north east);
    \fill[pattern={Dots[distance=4.4mm, radius=0.28mm]}, pattern color=plum!10!lavender!60!cream]
      (current page.south west) rectangle (current page.north east);
    \node at ([xshift=-3.2cm]current page.center) {\mascot[4.4cm]{happy}};
    \node[anchor=west, text=plum] at ([xshift=0.2cm, yshift=0.6cm]current page.center)
      {\titlefont\fontsize{40}{42}\selectfont Հարցեր};
    \node[anchor=west, text=plum, font=\mediumfont\large, align=left] at ([xshift=0.2cm, yshift=-1.0cm]current page.center)
      {Գրիր չաթում կամ\\ միացրու խոսափողը \emoji{megaphone}};
    \node[anchor=west, text=plum!70!cream, font=\small] at ([xshift=0.2cm, yshift=-2.1cm]current page.center)
      {Դին Բըրնեթ, «The Idiot Brain», 2016};
  \end{tikzpicture}
\end{frame}
